\section*{第 \S\arabic{exercisegroup} 节习题}
\phantomsection
\addcontentsline{toc}{section}{第 \protect\S\arabic{exercisegroup} 节习题}

\begin{enumerate}
\def\labelenumi{\arabic{enumi})}
\tightlist
\item
  设 I 是 R 中任意一个区间，A 是一个集合，g 是定义在 \(I \times A\) 上的有限实函数，使得对每个 \(\alpha \in A\) ，映射 \(t \mapsto g(t, \alpha)\) 在 I 上递减且 \(\geqslant 0\) ；再假设存在一个与 \(\alpha\) 无关的数 M，使得 \(g(t, \alpha) \leqslant M\) 在 \(I \times A\) 上成立 。
\end{enumerate}

\begin{enumerate}
\def\labelenumi{\alph{enumi})}
\item
  试证：若 \(\mathbf{f}\) 是 I 上的规则函数，且积分 \(\int_{\mathrm{I}} \mathbf{f}(t) dt\) 收敛，则积分 \(\int_{\mathrm{I}} \mathbf{f}(t) g(t, \alpha) dt\) 对 \(\alpha \in \mathrm{A}\) 一致收敛（利用第二中值定理；cf.~II, p.~82, exerc. 16）。
\item
  设 I = {[}a, +∞{]} ，并且当 t 趋于 +∞ 时 \(g(t, \alpha)\) （对 \(\alpha \in A\) ）一致地趋于 0。试证：若 f 是 I 上的规则函数，且存在数 k \textgreater{} 0 使得对含于 I 的每个紧区间 J 都有 \(\left\|\int_{J} \mathbf{f}(t) dt\right\| \leqslant k\) ，则积分 \(\int_{I} \mathbf{f}(t) g(t, \alpha) dt\) 对 \(\alpha \in A\) 一致收敛（同样的方法）。
\item
  设 \(I = [a, +\infty[\) 其中 a \textgreater{} 0，\(A = [0, +\infty[,\) 并且 \(g(t, \alpha) = \varphi(\alpha t)\) ，其中 \(\varphi\) 是 A 上的凸递减函数且 \(\geqslant 0\) ，当 x 趋于 \(+\infty\) 时趋于 0，并且满足 \(\varphi(0) = 1\) 。再设 \(f = h''\) ，其中 h 是 \([a, +\infty[,\) 上的二次可导函数，且它连同 \(h'\) 一起在该区间上有界。于是积分 \(\int_{a}^{+\infty} f(t) \varphi(\alpha t) dt\) 对 \(\alpha > 0\) 收敛；试证当 \(\alpha\) 趋于 0 时，它趋于 \(-\mathbf{h}'(a)\) （分部积分并利用第二中值定理）。
\end{enumerate}

* \(d\) ) 取 \(\varphi (x) = e^{-cx}\) ，其中 \(c > 0\) ；取 \(a = 1\) ，并且对 \(f\) 取复值函数 \(t\mapsto e^{i\log t} / t\) ，它是 I 上某个有界函数的导函数；试证：适当选取 \(c\) 后，积分 \(\int_{1}^{+\infty}e^{-\alpha ct}\frac{e^{i\log t}}{t} dt\) （它对 \(\alpha >0\) 绝对收敛）当 \(\alpha\) 趋于 0 时不趋于任何极限（分部积分之后作变量替换 \(\alpha t = u\) ；利用拉普拉斯变换理论可以看出，\(\int_0^{+\infty}e^{-cu + i\log u}du\) 对某些 \(c > 0\) 不为零。）*

\begin{enumerate}
\def\labelenumi{\arabic{enumi})}
\setcounter{enumi}{1}
\tightlist
\item
  设 \(f\) 与 \(g\) 是紧区间 \([a, b]\) 上的两个实规则函数，使得 \(f\) 在 \([a, b]\) 上递减，且 \(0 \leqslant g(t) \leqslant 1\) 。若令 \(\lambda = \int_{a}^{b} g(t) dt\) ，试证：
\end{enumerate}

\[
\int_ {b - \lambda} ^ {b} f (t) d t <   \int_ {a} ^ {b} f (t) g (t) d t <   \int_ {a} ^ {a + \lambda} f (t) d t
\]

除非 f 是常数，或者 g 在其连续的一切点处都等于 0（相应地，等于 1）（这时三项彼此相等）。（在积分 \(\int_{x}^{y} f(t)g(t)dt\) 中变动一个积分限。）

\begin{enumerate}
\def\labelenumi{\arabic{enumi})}
\setcounter{enumi}{2}
\tightlist
\item
  设 \(f(x, \alpha) = 1 / \sqrt{1 - 2\alpha x + \alpha^2}\) （ \(-1 < x < +1\) ，\(\alpha \in \mathbf{R}\) ）；试证函数 \(g(\alpha) = \int_{-1}^{+1} dx / \sqrt{1 - 2\alpha x + \alpha^2}\) 在 \(\mathbf{R}\) 上连续，但在 \(\alpha = 1\) 与 \(\alpha = -1\) 处没有导数；试证 \(f_{\alpha}'(x, \alpha)\) 对一切 \(\alpha \in \mathbf{R}\) 以及对 \(x \in I = ] - 1, +1[\) 都存在，且在 \(I \times R\) 上连续，并且积分 \(\int_{-1}^{+1} f_{\alpha}'(x, \alpha) dx\) 对一切 \(\alpha \in R\) 都存在，但要注意该积分在点 \(\alpha = 1\) 或点 \(\alpha = -1\) 的邻域内不是一致收敛的。
\end{enumerate}

¶ 4) 设 I 是 \(\mathbf{R}\) 中一个区间，A 是点 \(\alpha_0\) 在域 \(\mathbf{R}\) （相应地，域 \(\mathbf{C}\) ）中的一个邻域，\(\mathbf{f}\) 是 \(\mathrm{I} \times \mathrm{A}\) 到完备赋范空间 \(\mathrm{E}\) （\(\mathbf{R}\) 上的）中的连续映射，使得 \(\mathbf{f}_{\alpha}'(x, \alpha)\) 存在且在 \(\mathrm{I} \times \mathrm{A}\) 上连续。设 \(a(\alpha)\) ，\(b(\alpha)\) 是定义在 A 上、取值于 I 中的两个连续函数，使得在 A 上恒有 \(\mathbf{f}(a(\alpha), \alpha) = \mathbf{f}(b(\alpha), \alpha) = 0\) 。试证：函数 \(\mathbf{g}(\alpha) = \int_{a(\alpha)}^{b(\alpha)} \mathbf{f}(t, \alpha) dt\) 有导数，等于 \(\int_{a(\alpha_0)}^{b(\alpha_0)} \mathbf{f}_{\alpha}'(t, \alpha_0) dt\) （在点 \(\alpha_0\) 处），即使 \(a\) 与 \(b\) 在点 \(\alpha_0\) 处不可导（设 M 是 \(\left\| \mathbf{f}_{\alpha}'(x, \alpha) \right\|\) 在 \((b(\alpha_0), \alpha_0)\) 的一个紧邻域上的上确界；注意，对 \(b(\alpha)\) 应用博尔扎诺定理可知，对每一点 \(x\) ，只要它属于以 \(b(\alpha_0)\) 与 \(b(\alpha)\) 为端点的区间，且 \(\alpha\) 充分接近 \(\alpha_0\) ，便有 \(\| \mathbf{f}(x, \alpha) \| \leqslant M |\alpha - \alpha_0|\) ）。

\begin{enumerate}
\def\labelenumi{\arabic{enumi})}
\setcounter{enumi}{4}
\tightlist
\item
  设 \(\mathbf{f}\) 是紧区间 \(I = [0, a]\) 上的连续向量函数。试证：若在点 \(\alpha_0 \in I\) 处存在 \(\varepsilon > 0\) ，使得 \((\mathbf{f}(x) - \mathbf{f}(\alpha_0)) / |x - \alpha_0|^{1/2 + \varepsilon}\) 当 \(x\) 趋于 \(\alpha_0\) 时保持有界，则函数 \(\mathbf{g}(\alpha) = \int_0^\alpha \mathbf{f}(x) dx / \sqrt{\alpha - x}\) 有导数，等于
\end{enumerate}

\[
\frac {1}{\sqrt {\alpha_ {0}}} \mathbf {f} (\alpha_ {0}) - \frac {1}{2} \int_ {0} ^ {\alpha_ {0}} \frac {\mathbf {f} (x) - \mathbf {f} (\alpha_ {0})}{(\alpha_ {0} - x) ^ {3 / 2}} d x
\]

即：在点 \(\alpha_0\) 处（当 \(\alpha_0 > 0\) 时）取上值为导数，而当 \(\alpha = 0\) 时导数为 0 。

当 \(\mathbf{f}\) 是实函数 \(\sqrt{\alpha_0 - x}\) 时，试证函数 \(\mathbf{g}\) 在点 \(\alpha_0\) 处有无穷导数。

¶ 6) 设 \(\mathrm{I} = [a, b]\) ，\(\mathrm{A} = [c, d]\) 是 \(\mathbf{R}\) 上的两个紧区间；设 \(\mathbf{f}\) 是定义在 \(\mathrm{I} \times \mathrm{A}\) 上、取值于完备赋范空间 \(\mathrm{E}\) （\(\mathbf{R}\) 上的）中的函数，使得：对一切 \(\alpha \in \mathrm{A}\) ，映射 \(t \mapsto \mathbf{f}(t, \alpha)\) 在 \(\mathrm{I}\) 上是规则的；\(\mathbf{f}\) 在 \(\mathrm{I} \times \mathrm{A}\) 上有界；并且集 \(\mathrm{D}\) ——即 \(\mathbf{f}\) 在 \(\mathrm{I} \times \mathrm{A}\) 中的不连续点所成的集——与每条直线 \(x = x_0\) 以及每条直线 \(\alpha = \alpha_0\) 都只交于有限多个点（ \(x_0 \in \mathrm{I}\) ，\(\alpha_0 \in \mathrm{A}\) ）。

\begin{enumerate}
\def\labelenumi{\alph{enumi})}
\item
  试证函数 \(\mathbf{g}(\alpha) = \int_{a}^{b} \mathbf{f}(t, \alpha) dt\) 在 A 上连续（给定 \(\alpha_0 \in \mathrm{A}\) 与 \(\varepsilon > 0\) ，证明存在 \(\alpha_0\) 的一个邻域 V 以及含于 I 中的有限多个区间 \(J_k\) ，其长度之和 \(\leqslant \varepsilon\) ，使得若用 J 表示 \(\bigcup_{k} J_k\) 在 I 中的余集，则 \(\mathbf{f}\) 在 \(\mathrm{J} \times \mathrm{V}\) 上连续）。
\item
  试证积分次序交换公式（II, p.~77，公式 (15)）仍然成立（用与 a) 中相同的方法）。
\end{enumerate}

\begin{enumerate}
\def\labelenumi{\arabic{enumi})}
\setcounter{enumi}{6}
\tightlist
\item
  设 \(f\) 是在区间 \(]0, +\infty[\) 上定义且具有连续导数的实函数，并且使得
\end{enumerate}

\[
\lim _ {x \to 0} f (x) = \lim _ {x \to + \infty} f (x) = 0.
\]

积分 \(\int_0^{+\infty}f'(\alpha t)dt\) 在每个有界区间 \(]0,a]\) 上都有定义且连续；试证积分 \(\int_0^{+\infty}dt\int_0^a f'(\alpha t)d\alpha\) 可能不存在，也可能不同于

\[
\int_ {0} ^ {a} d \alpha \int_ {0} ^ {+ \infty} f ^ {\prime} (\alpha t) d t.
\]

¶ 8) 设 \(I\) 与 \(J\) 是 \(\mathbf{R}\) 中两个任意区间，\(\mathbf{f}\) 是定义在 \(I \times J\) 上且连续、取值于完备赋范空间 \(E\) （\(\mathbf{R}\) 上的）中的函数。假设：

\(1^{\circ}\) 积分 \(\int_{\mathrm{I}}\mathbf{f}(x,y)dx\) 当 \(y\) 跑遍含于 J 的任一紧区间时一致收敛；

\(2^{\circ}\) 积分 \(\int_{\mathrm{J}}\mathbf{f}(x,y)dy\) 当 \(x\) 跑遍含于 I 的任一紧区间时一致收敛；

\(3^{\circ}\) 若对含于 I 的每个紧区间 \(\mathrm{H}\) 令 \(\mathbf{u}_{\mathrm{H}}(y) = \int_{\mathrm{H}} \mathbf{f}(x, y) dx\) ，则积分 \(\int_{\mathrm{J}} \mathbf{u}_{\mathrm{H}}(y) dy\) 对 \(\mathrm{H} \in \mathfrak{K}(\mathrm{I})\) 一致收敛（此处取的是含于 I 的紧区间所成的右向有序集）。

在这些条件下，试证积分 \(\int_{\mathrm{I}} dx \int_{\mathrm{J}} \mathbf{f}(x, y) dy\) 与 \(\int_{\mathrm{J}} dy \int_{\mathrm{I}} \mathbf{f}(x, y) dx\) 都存在且相等。

* 9) 由习题 8 推出，积分

\[
\int_ {0} ^ {\infty} d x \int_ {0} ^ {\infty} e ^ {- y x ^ {2}} \sin y d y \quad \text { and } \quad \int_ {0} ^ {\infty} d y \int_ {0} ^ {\infty} e ^ {- y x ^ {2}} \sin y d x
\]

存在且相等。*

\begin{enumerate}
\def\labelenumi{\arabic{enumi})}
\setcounter{enumi}{9}
\tightlist
\item
  \begin{enumerate}
  \def\labelenumii{\alph{enumii})}
  \tightlist
  \item
    若 \(h, u, v'\) 是开区间 \(]a, b[\) （\(\mathbf{R}\) 中的）上实规则函数的原函数，若 \(v\) 是 \(v'\) 的原函数，并且 \(v(x) > 0\) 在该区间上成立，则有雷德赫费尔恒等式
  \end{enumerate}
\end{enumerate}

\[
h u ^ {\prime 2} = - h \mathrm{D} \left(\frac {h v ^ {\prime}}{v}\right) + h v ^ {2} \left(\mathrm{D} \left(\frac {u}{v}\right)\right) ^ {2} + \mathrm{D} \left(\frac {u ^ {2} h v ^ {\prime}}{v}\right)
\]

在 {]}a, b{[} 中使诸导数有定义的点处成立。

* b) 设 \(v\) ，\(w\) 是两个 \(> 0\) 的函数，定义在 \(]a\) ，\(b[\) 上，它们分别是规则函数 \(v' > 0\) 与 \(w' \leqslant 0\) 的原函数。由 \(a)\) 推出：对每个函数 \(u\) ，只要它是规则函数 \(u'\) 在 \(]a\) ，\(b[\) 上的原函数且满足 \(\liminf_{x \to a, x > a} u(x) = 0\) ，则由积分 \(\int_{a}^{b} \frac{w(x)}{v'(x)} u'^2(x) dx\) 收敛这一假设可以推出积分 \(\int_{a}^{b} \frac{w'(x)}{v(x)} u^2(x) dx\) 收敛，并且

\[
\limsup _ {x \to b, x <   b} u ^ {2} (x) w (x) / v (x)
\]

是有限的，而且有不等式

\[
\int_ {a} ^ {b} \frac {w (x)}{v ^ {\prime} (x)} u ^ {\prime 2} (x) d x \geqslant - \int_ {a} ^ {b} \frac {w ^ {\prime} (x)}{v (x)} u ^ {2} (x) d x + \operatorname * {l i m s u p} _ {x \to b, x <   b} \frac {u ^ {2} (x) w (x)}{v (x)}.\tag{*}
\]

（令 \(h = w / v'\) ，首先注意 \(\int_{c}^{x} dt / h(t) \leqslant v(x) / w(x)\) 对 \(a < c < x < b\) 成立，并注意由柯西–施瓦茨不等式

\[
(u (x) - u (c)) ^ {2} \leqslant \left(\int_ {c} ^ {x} h (t) u ^ {\prime 2} (t) d t\right) \left(\int_ {c} ^ {x} \frac {d t}{h (t)}\right),\tag{**}
\]

可以推出

\[
\liminf _ {x \to a, x > a} u ^ {2} (x) w (x) / v (x) = 0,
\]

然后对雷德赫费尔恒等式积分。）

\begin{enumerate}
\def\labelenumi{\alph{enumi})}
\setcounter{enumi}{2}
\tightlist
\item
  由 (*) 推出：若 u 是 {]}0, 1{]} 上一个规则函数的原函数，若
\end{enumerate}

\[
\liminf _ {x \to 0, x > 0} u (x) = 0
\]

且积分 \(\int_0^1 u'^2 (t)dt\) 收敛，则 \(\int_0^1 (u(t) / t)^2 dt\) 也收敛，并且对一切 \(\alpha > 0\) 有不等式

\[
\int_ {0} ^ {1} u ^ {\prime 2} (t) d t \geqslant \alpha (1 - \alpha) \int_ {0} ^ {1} \left(\frac {u (t)}{t}\right) ^ {2} d t + \alpha u ^ {2} (1).
\]

\begin{enumerate}
\def\labelenumi{\alph{enumi})}
\setcounter{enumi}{3}
\tightlist
\item
  设 \(\alpha\) 是一个 \textgreater0 的数，K 是一个 \textgreater0 、可导且递减的函数，定义在 \(]0, +\infty[\) 上，使得 \(\lim_{x \to +\infty} K(x) = 0\) 。若 u 是 \(]0, +\infty[,\) 上一个规则函数的原函数，满足 \(\liminf_{x \to 0, x > 0} u(x) = 0\) ，并且积分 \(\int_{0}^{+\infty} x^{1-\alpha} K(x) u'^{2}(x) dx\) 收敛，则函数 \(x^{-\alpha} K(x) u^{2}(x)\) 当 x 趋于 0 或趋于 \(+\infty\) 时趋于 0 ，且积分
\end{enumerate}

\[
\int_ {0} ^ {+ \infty} x ^ {- \alpha} \mathrm{K} ^ {\prime} (x) u ^ {2} (x) d x
\]

收敛，并且有

\[
\int_ {0} ^ {+ \infty} x ^ {1 - \alpha} \mathrm{K} (x) u ^ {\prime 2} (x) d x \geqslant - \alpha \int_ {0} ^ {+ \infty} x ^ {- \alpha} \mathrm{K} ^ {\prime} (x) u ^ {2} (x) d x
\]

（在雷德赫费尔恒等式中取 \(v(x) = x^{\alpha}\) ，\(h(x) = x^{1 - \alpha}\mathrm{K}(x)\) ）。

特别地，取 \(\alpha = \frac{1}{2}\) 与 \(\mathrm{K}(x) = x^{-1 / 2}\) ，有

\[
4 \int_ {0} ^ {+ \infty} u ^ {\prime 2} (x) d x \geqslant \int_ {0} ^ {+ \infty} \left(\frac {u (x)}{x}\right) ^ {2} d x
\]

（哈代–李特尔伍德不等式）。

\begin{enumerate}
\def\labelenumi{\alph{enumi})}
\setcounter{enumi}{4}
\tightlist
\item
  设 \(\alpha \geqslant -1\) ，设 \(K\) 是一个 \(\geqslant 0\) 、可导且递增的函数，定义在 \([0, +\infty[\) 上。假设积分
\end{enumerate}

\[
\int_ {0} ^ {+ \infty} x ^ {- \alpha} \mathrm{K} (x) u ^ {\prime 2} (x) d x \quad \text { and } \quad \int_ {0} ^ {+ \infty} x ^ {\alpha} \mathrm{K} (x) u ^ {2} (x) d x
\]

收敛。则 \(\mathrm{K}(x)u^{2}(x)\) 当 x 趋于 \(+\infty\) 时趋于 0 ，并且有

\[
\int_ {0} ^ {+ \infty} \mathrm{K} ^ {\prime} (x) u ^ {2} (x) d x \leqslant 2 \left(\int_ {0} ^ {+ \infty} x ^ {- \alpha} \mathrm{K} (x) u ^ {\prime 2} (x) d x\right) ^ {1 / 2} \left(\int_ {0} ^ {+ \infty} x ^ {\alpha} \mathrm{K} (x) u ^ {2} (x) d x\right) ^ {1 / 2}
\]

（在雷德赫费尔恒等式中取 \(v(x) = \exp(-cx^{\alpha})\) ，其中取适当的常数 \(c\) ）。

特别地，若 a 与 b 是使 \(b + 1 \geqslant a\) 且 \(a + b \geqslant 0\) 的常数，则有

\[
(a + b) \int_ {0} ^ {+ \infty} x ^ {a + b - 1} u ^ {2} (x) d x \leqslant 2 \left(\int_ {0} ^ {+ \infty} x ^ {2 a} u ^ {\prime 2} (x) d x\right) ^ {1 / 2} \left(\int_ {0} ^ {+ \infty} x ^ {2 b} u ^ {2} (x) d x\right) ^ {1 / 2}
\]

只要右端的积分收敛（推广的 H. 外尔不等式）。

\(f)\) 若 \(0 < \alpha < 2\) ，且 \(u\) 是 \([0, \alpha]\) 上一个规则函数的原函数，并满足 \(u(0) = 0\) ，则有

\[
\int_ {0} ^ {\alpha} (\alpha - x) ^ {\alpha - 1} x ^ {1 - \alpha} u ^ {\prime} (x) \left(u ^ {\prime} (x) - 2 u (x)\right) d x \geqslant 0
\]

（推广的奥皮亚尔不等式）。（适当应用 (*) ，把 \(u(x)\) 换成 \(e^{-x}u(x)\) 。）

\(g\) ) 若 \(u\) 是 \([0,b]\) 上一个规则函数的原函数，且 \(u(0) = 0\) ，则

\[
\int_ {0} ^ {b} e ^ {- 2 x} u ^ {\prime 2} (x) d x \geqslant \int_ {0} ^ {b} e ^ {- 2 x} u (x) u ^ {\prime} (x) d x + \frac {1}{2} u ^ {2} (b) e ^ {- 2 b}
\]

（赫拉夫卡不等式）（用与 \(f\) 中相同的方法）。

\begin{enumerate}
\def\labelenumi{\alph{enumi})}
\setcounter{enumi}{7}
\tightlist
\item
  若 u 是 R 上一个规则函数的原函数，则对一切 \(t \in R\) ，
\end{enumerate}

\[
| u (t) | \leqslant \left(\int_ {- \infty} ^ {+ \infty} u ^ {\prime 2} (x) d x\right) ^ {1 / 4} \left(\int_ {- \infty} ^ {+ \infty} u ^ {2} (x) d x\right) ^ {1 / 4}
\]

只要右端的两个积分收敛（考虑两个区间 \(]-\infty,t]\) 与 \([t,+\infty[\) ；在这两个区间上都取 \(v(x)=e^{\alpha x}\) ，在第一个区间上取 \(h(x)=1/\alpha\) ，在第二个区间上取 \(h(x)=-1/\alpha\) ，然后适当选取 \(\alpha>0\) ）。\(^{*}\)

